\section{Planteamiento del problema}

El polimetilmetacrilato (PMMA), o acrílico, es un termoplástico amorfo con transición vítrea $T_g \approx$ \SI{105}{\celsius} \cite{gilormini2010,brandrup1999} y ventana de conformado por línea térmica entre \num{150} y \SI{160}{\celsius} \cite{plexiglas_forming,throne2008}. Su transparencia y facilidad de mecanizado lo hacen el material estructural más usado en montajes experimentales y prototipos docentes.

El doblado por línea (\emph{line bending}) calienta una franja estrecha por encima de $T_g$ hasta el estado gomoso, la pliega y sostiene el ángulo mientras enfría \cite{osswald2012,throne2008}. La calidad depende de tres variables acopladas: el perfil de temperatura a lo ancho de la franja, que fija la zona plástica; el tiempo de permanencia, que gobierna la penetración térmica y la relajación de esfuerzos; y el enfriamiento bajo carga, que determina la recuperación elástica (\emph{springback}). Calentar de menos produce microfisuración (\emph{crazing}) y springback alto; calentar de más genera burbujas y deformación fuera de la línea \cite{yan2023,plexiglas_forming}. Enfriar bajo restricción relaja los esfuerzos internos \cite{osswald2012,yan2023}.

La oferta comercial se reparte entre dos extremos: \emph{strip heaters} portátiles sin sensor ni realimentación (\SIrange{300}{1500}{\watt}, 80 a 300 USD), donde el operario ajusta por tiempo y pliega a pulso contra plantilla \cite{shannon_line}, y dobladoras de control numérico con prensado motorizado, pero de \SIrange{1200}{3000}{\milli\meter} de longitud útil y precios de miles de dólares \cite{shannon_bending,mic_auto}. No se identificó oferta intermedia: equipos de mesa, de longitud útil inferior a \SI{300}{\milli\meter}, con control realimentado y plegado automático, que es justo lo que necesita un laboratorio universitario. En el Laboratorio de Mecatrónica de EAFIT el doblado se hace hoy a mano sobre resistencias sin instrumentar: el resultado depende de la destreza del técnico y no es reproducible.

\noindent\textbf{Pregunta de investigación:} \emph{¿Puede una máquina de doblado por línea térmica de pequeño formato y bajo costo, con control realimentado de temperatura y plegado motorizado, producir dobleces en PMMA con error angular no mayor a \SI{\pm 2}{\degree} y error de posición no mayor a \SI{\pm 1}{\milli\meter}?}
